\chapter{Estado da arte}
\label{chap:arte}

\section{Historia da Gameboy}
O 21 de abril de 1989, a empresa de videoxogos Nintendo lanzou no mercado nipón o modelo orixinal da Game Boy, 
con nome en clave DMG-01, baixo a direción de Yunpei Yokoi e Satoru Okada cun prezo de saída de 12\,500 ¥.
Con este modelo Nintendo entrou no mercado dos videoxogos portátiles con cartuchos intercambiables.
Xunto ca consola, nos mercados estadounidense e europea lanzouse xunto co xogo Tetris,
o que contribuíu a que tanto este xogo como a consola se convertesen en grandes éxitos,
con 35 e 118 millóns de unidades vendidas, xa que axudou a popularizala nunha franxa de mercado máis adulta % todo: reformular
Tamén foi a consola que dou vida a outro fenómeno de masas,
a primeira xeración de Pokémon, que xuntando todas as súas versións vendeu 46 millóns de unidades a nivel mundial.

\subsection{Revisions}

Ao longo de máis de quince anos, Nintendo lanzou ata sete revisións de hardware
dentro da familia Game Boy, que se resumen na táboa % ~\ref{tab:gb-modelos}
% (páxina~\pageref{tab:gb-modelos}). 
Os modelos divídense en dúas xeracións
compatibles entre si en canto a software: a orixinal (DMG, Pocket, Light e
Color, todos baseados no procesador SM83/LR35902) e a Game Boy Advance (GBA,
SP e Micro, baseados no ARM7TDMI cun coprocesador Z80 para compatibilidade
cara atrás). Os prezos refírense ao lanzamento no mercado xaponés, primeiro
mercado de todas estas consolas.

Baixo a marca Game Boy atopamos dúas familias de consolas. A Game Boy orixinal, contou con tres revisións,
a Game Boy Pocket, cun formato reducido; a Game Boy Light (que soamente saliu en Xapón); con retroiluminación
e a Game Boy Color, con unha pantalla a cores e un procesador cun modo de dobre velocidade.

Despois do éxito da Game Boy Orixinal, Nintendo continuo a súa aposta polas videoconsolas portátiles ca
Game Boy advance, que conta con tres modelos, a Game Boy Advance,
a Game Boy Advance SP, cun formato pregable, pantalla retroiluminada e con batería recargable en vez de pilas,
e a Game Boy Advance Micro, cun formato reducido.

Ademais destas consolas portátiles, Nintendo lanzou en Xapón a Super Game Boy e Super Game Boy 2.
A súa función era permitir a execución de xogos das portátiles nun televisor
Estes periféricos para a Super Famicom incluían todo o hardware dunha Game Boy,
unicamente cambiando a saída de vídeo e audio hacia un televisor.


Os dous primeiros modelos contaban con retrocompatibilidade ca familia orixinal, incluindo o hardware dunha 
Game Boy Color. O modelo Micro en cambio non contaba con esta retrocompatibilidade.
A partir de aquí as seguintes consolas, dende a Nintendo DS e Nintendo DS Lite,
con retrocompatibilidade soportada oficialmente e lector de cartuchos,
ata o último modelo da familia Nintendo 3DS, soamente contaban con compatibilidade ca familia Advance,
empregando emulación para executar os xogos da familia orixinal.


% TODO: axustar frecuencias e facer distincion mebihercios e megahercios e revisar datos
\begin{table}[hp!]
  \centering
  \resizebox{\textwidth}{!}{%
  \rowcolors{2}{white}{udcgray!25}%
  \begin{tabular}{l|c|c|l|l|l|r}
  \rowcolor{udcpink!25}
  \textbf{Modelo} & \textbf{Lanz.} & \textbf{Fin fab.} & \textbf{Procesador} & \textbf{Pantalla} & \textbf{Alimentación} & \textbf{Prezo} \\\hline
  \textit{Game Boy} (DMG-01)         & 1989 & 2003 & Sharp SM83 \mbox{4\,194\,304\,Hz}              & STN LCD mono.        & 4 $\times$ AA       & 12\,500\,¥ \\
  \textit{Game Boy Pocket} (MGB)     & 1996 & 2003 & Sharp SM83 \mbox{4\,194\,304\,Hz}              & FSTN LCD mono.       & 2 $\times$ AAA      & 6\,800\,¥  \\
  \textit{Game Boy Light} (MGB-101)  & 1998 & 2003 & Sharp SM83 \mbox{4\,194\,304\,Hz}              & FSTN retroil.        & 2 $\times$ AA       & 6\,800\,¥  \\
  \textit{Game Boy Color} (CGB)      & 1998 & 2003 & Sharp SM83 ata \mbox{4\,194\,304/8\,388\,608\,Hz}       & TFT LCD a cor        & 2 $\times$ AA       & 8\,900\,¥  \\
  \textit{Game Boy Advance} (AGB)    & 2001 & 2009 & ARM7TDMI \mbox{4\,194\,304\,Hz} MHz + Sharp SM83                              & TFT cor panor.       & 2 $\times$ AA       & 9\,800\,¥  \\
  \textit{Game Boy Advance SP} (AGS) & 2003 & 2009 & ARM7TDMI + Sharp SM83                                & TFT pregable         & Li-ion recargable   & 12\,500\,¥ \\
  \textit{Game Boy Micro} (OXY)      & 2005 & 2009 & ARM7TDMI                                & TFT retroiluminada   & Li-ion recargable   & 12\,000\,¥ \\
  \end{tabular}}
  \caption{Familia Game Boy: revisión, especificacións principais e datas comerciais. Prezos de saída no mercado xaponés.}
  % \label{tab:gb-modelos}
\end{table}

\subsection{Hardware como producto}
Esta consola conta ca configuración básica de botóns de Nintendo herdado dos Game \& Watch e da NES.
Contaba cunha pantalla LCD STN reflectiva monocroma de 2,5 polgadas,
160$\times$144 píxeles e catro tons de gris (codificados con dous bits), sen retroiluminación.
Para alimentarse empregaba 4 pilas AA que acaban un autonomía de entre 15 e 30 horas, segundo

Estes factores, fana parecer unha consola modesta para a súa época con tecnoloxía anticuada, 
sobretodo comparandoa cos seus principales competidores.
Comparada cos seus rivais directos, a Game Boy era a opción tecnoloxicamente máis modesta.
A Atari Lynx (1989, 179,95 \$ USD ao lanzamento) ofrecía pantalla en cor retroiluminada,
escalado e rotación de sprites por hardware e unha paleta de 4\,096 cores, 
pero o seu consumo dunha sextupla AA reducía a autonomía a tres ou catro horas  % \cite{wiki_lynx}.
A Sega Game Gear (1990 en Xapón, 19\,800 ¥; 149,99 \$ USD en EUA) baseábase nun Z80 emparentado co
Master System  e podía mostrar 32 cores simultáneas,
mais sufría idénticos problemas de duración (3-5 horas con seis pilas AA)
e un prezo notablemente superior ao da Game Boy %  \cite{wiki_gamegear}.

Estes factores foron decisivos para o seu enorme éxito. Nintendo optou por ter unha consola
realmente portátil usando tecnoloxía modesta e barata para o consumidor.
Fronte aos e 10,62 millóns da Game Gear e os menos de 3 millóns da Lynx, a Game Boy suporeunas
notablemente grazas á súa estratexia de posicionamento no mercado portatil % cita


\section{Escena da emulación de GameBoy}

\subsection{Emulación en xeral}
A emulación é a técnica que permite executar software deseñado para unha plataforma nunha plataforma diferente,
replicando mediante software o comportamento do hardware orixinal: o procesador
, a memoria e dos periféricos da máquina orixinal de xeito que se consigue un funcionamento equivalente ao da máquina orixinal.

Non todos os emuladores abordan este problema co mesmo alcance.
Na emulación completa, ou emulación de sistema, replícase a totalidade do hardware da máquina obxectivo,
dende o procesador ata o hardware gráfico e de son.
É a aproximación habitual na emulación de consolas, e tamén a que segue este traballo.
Noutros contextos abonda cunha emulación parcial.

O caso máis representativo é QEMU % \cite{qemu}
, que ofrece os dous alcances: no seu modo de sistema emula unha máquina completa,
capaz de arrancar un sistema operativo enteiro,
mentres que no seu modo de usuario executa binarios doutra arquitectura
sobre o kernel nativo, emulando unicamente o espazo de usuario.

Esta segunda aproximación é a que segue FEX~\cite{fexemu},
que permite executar binarios x86 e x86-64 sobre sistemas GNU/Linux con procesadores ARM64:
as chamadas ao sistema trasládanse directamente ao kernel nativo,
e as librarías gráficas poden substituírse por implementacións nativas,
evitando así ter que emular hardware ningún.
Unha solución similar é Rosetta 2~\cite{rosetta2},
desenvolvida por Apple para facilitar a transición de macOS aos seus procesadores ARM64,
que traduce as aplicacións x86-64 a código nativo no momento da súa instalación,
recorrendo á tradución dinámica só cando é imprescindible.

En canto á execución do código emulado, a aproximación máis sinxela é a interpretación,
na que o emulador decodifica e executa as instrucións do programa unha a unha.
Para plataformas máis potentes esta técnica resulta insuficiente,
e recórrese á compilación \textit{just-in-time} (JIT)~\cite{jit},
que traduce bloques de código da máquina emulada a código nativo durante a execución.
Os bloques traducidos gárdanse nunha caché de código,
de xeito que cada bloque só se traduce a primeira vez que se executa
e as seguintes execucións reutilizan directamente a versión nativa,
amortizando así o custo da tradución no código que se executa con máis frecuencia.

Algo semellante sucede co hardware gráfico.
Nas consolas máis sinxelas este emúlase por software,
empregando a GPU do sistema anfitrión unicamente para presentar a imaxe final,
pero nas máquinas con hardware 3D as ordes de debuxado do xogo
tradúcense a chamadas de APIs gráficas modernas,
sendo a GPU real quen realiza o traballo de renderización.
Esta tradución esixe xerar e compilar durante a execución shaders equivalentes
a cada configuración do pipeline gráfico orixinal,
un proceso custoso que provoca paróns perceptibles durante o xogo,
polo que os emuladores manteñen tamén unha caché de shaders,
persistida en disco entre sesións, para que cada shader só se compile unha vez.
Un exemplo representativo destas técnicas é Dolphin~\cite{Dolphin},
emulador de GameCube e Wii, que combina un compilador JIT de PowerPC a código nativo
coa emulación do pipeline gráfico da consola sobre a GPU.
Ademais da caché de shaders, este emulador introduciu os \textit{ubershaders}~\cite{DolphinUbershaders},
shaders xenéricos que interpretan a configuración do pipeline na propia GPU,
eliminando os paróns de compilación mesmo na primeira execución.

Existen ademais alternativas á emulación que parten do mesmo labor de enxeñería inversa.
A descompilación busca reconstruír o código fonte orixinal do xogo a partir do seu binario,
o que permite despois compilalo de forma nativa para outras plataformas;
proxectos como Ship of Harkinian~\cite{ShipofHarkinian}
, port nativo de The Legend of Zelda: Ocarina of Time,
ou Dusklight~\cite{Dusklight}
nacen de descompilacións deste type.
A recompilación estática, pola súa banda, traduce o binario orixinal unha única vez,
antes da execución, a código para a plataforma de destino,
sen necesidade de reconstruír o código fonte nin de manter un emulador en tempo de execución.
Ferramentas como N64Recomp~\cite{N64Recomp}
, para Nintendo 64, ou XenonRecomp~\cite{XenonRecomp}
, para Xbox 360, permitiron a creación de ports nativos de xogos comerciais con melloras
imposibles nun emulador tradicional, como maiores resolucións e taxas de refresco.
É o caso de Zelda64Recomp~\cite{Zelda64Recomp}
, construído sobre N64Recomp para portar The Legend of Zelda: Ocarina of Time,
e de UnleashedRecomp~\cite{UnleashedRecomp}
, construído sobre XenonRecomp para portar Sonic Unleashed.

\subsection{Historia emuladores de Game Boy}

A emulación de Game Boy susténtase nun enorme labor de enxeñería inversa e documentación
realizado pola comunidade ao longo de tres décadas, xa que Nintendo nunca publicou
a documentación técnica oficial da consola.
O documento de referencia son os Pan Docs % \cite{pandocs}
, orixinados a mediados dos anos noventa e mantidos na actualidade pola comunidade gbdev,
que describen o comportamento de todos os compoñentes do sistema.
Complementándoos, o Game Boy: Complete Technical Reference % \cite{gbctr}
de Joonas Javanainen documenta o hardware con precisión de ciclo,
a partir de medicións sobre consolas reais.
Este coñecemento acumulado deu lugar tamén a un ecosistema de desenvolvemento activo:
RGBDS % \cite{rgbds}
ofrece un ensamblador e enlazador completos para a creación de novos xogos,
e GB Studio % \cite{gbstudio}
permite desenvolvelos visualmente sen necesidade de programar,
o que mantén viva a escena homebrew da consola.

Os primeiros emuladores apareceron a mediados dos anos noventa,
coa consola aínda á venda.
Entre os pioneiros destacan Virtual GameBoy de Marat Fayzullin
e NO\$GMB de Martin Korth, que xa incluía un completo depurador
e cuxo autor foi ademais responsable da primeira versión dos Pan Docs.
Na actualidade os emuladores máis relevantes son proxectos de código aberto
cun forte foco na precisión:
SameBoy % \cite{sameboy}
, considerado un dos máis fieis ao hardware orixinal,
e mGBA % \cite{mgba}
, centrado na Game Boy Advance pero con soporte da familia orixinal,
son os máis empregados, tanto de forma independente como integrados en RetroArch % \cite{retroarch}
, un frontend multiplataforma que agrupa ducias de emuladores como núcleos intercambiables.
Mención aparte merece Mooneye GB % \cite{mooneye}
, tamén de Joonas Javanainen, desenvolvido como ferramenta de investigación do hardware
e acompañado dunha suite de tests que se converteu nunha referencia
para validar a precisión doutros emuladores.

As técnicas máis avanzadas descritas na sección anterior tiveron tamén
os seus primeiros experimentos nesta plataforma,
como gb-recompiled % \cite{gbrecompiled}
, unha proba de concepto de recompilación estática de xogos de Game Boy.
Porén, esta aproximación resulta aquí máis complicada que noutras plataformas:
os xogos de Game Boy executan con frecuencia código dende a RAM
e empregan código automodificable,
o que impide traducir estaticamente todo o programa
e obriga a manter un intérprete de apoio en tempo de execución.

\section{Uso comercial de emuladores de Game Boy}

A propia Nintendo emprega a emulación para explotar comercialmente o seu catálogo clásico.
A Consola Virtual da Nintendo 3DS ofrecía xogos de Game Boy e Game Boy Color
executados sobre un emulador propio,
mentres que os xogos de Game Boy Advance distribuídos nesa consola
se executaban de forma nativa,
aproveitando que o hardware da 3DS inclúe o das súas predecesoras Nintendo DS e Game Boy Advance.
Na actualidade, o servizo Nintendo Switch Online inclúe catálogos de xogos
de Game Boy e Game Boy Advance executados sobre emuladores desenvolvidos
pola súa filial europea Nintendo European Research \& Development (NERD). % \cite{nso}

\subsection{Consolas de emulación}
Nos últimos anos consolidouse ademais un mercado de consolas portátiles
dedicadas exclusivamente á emulación,
con fabricantes como Anbernic, Miyoo ou Powkiddy % \cite{anbernic}
que lanzan constantemente novos modelos cun formato que imita o das portátiles clásicas.
Trátase de dispositivos con procesadores ARM e sistemas GNU/Linux ou Android
que executan frontends como RetroArch, EmulationStation, iiSU ou CocoonFE;
e cuxa principal utilidade é xogar comodamente a catálogos retro de múltiples plataformas.
Dende o punto de vista da emulación de Game Boy non proporcionan ningunha vantaxe técnica, xa que
o seu hardware é ordes de magnitude máis potente do necesario para emular esta consola,
sendo capaces de emular sistemas máis complexos.

\subsection{Emulación baseada en hardware}
Unha aproximación radicalmente distinta á emulación por software
é a reimplementación do hardware orixinal sobre unha \acrfull{fpga},
un circuíto integrado reconfigurable que permite describir nunha \acrfull{hdl}
o comportamento dos chips orixinais e executalo directamente en silicio.
Fronte a un emulador por software, que reproduce o comportamento do sistema
de forma secuencial, unha \acrshort{fpga} replica o paralelismo real do hardware,
o que permite acadar unha precisión e unha latencia moi difíciles de igualar por software.

O exemplo máis destacado no ámbito da Game Boy é a Analogue Pocket % \cite{analoguepocket}
, unha consola portátil que reimplementa en \acrshort{fpga} o hardware das familias
Game Boy e Game Boy Advance, entre outras portátiles clásicas,
e que é compatible cos cartuchos e accesorios orixinais.
A súa plataforma aberta, openFPGA, permite ademais á comunidade desenvolver
núcleos para outros sistemas.
A cambio, trátase dunha solución notablemente máis cara que calquera emulador por software,
cun prezo de saída de 239,99 \$ USD,
dirixida a un público que prioriza a fidelidade ao hardware orixinal fronte unha solución económica
que permita xogar ao mesmo catálogo de xogos.
